\chapter{Adição com reagrupamento}\label{ch:g2:addition}

Quando as unidades de uma \omterm{def:g1:addition:def}{soma} passam de nove, dez delas se agrupam
numa dezena nova --- e esse pacotinho, levado para a coluna
seguinte, é o famoso \emph{vai-um}. Este capítulo explica o vai-um
com palitos e depois treina a conta armada.

\section{Por que existe o vai-um}

\begin{example}[Agrupar dez unidades]\label{ex:g2:addition:bundle}
$27 + 15$: junte as unidades, $7 + 5 = 12$ unidades. Mas dez
unidades formam uma dezena! Amarre dez delas: sobram $2$ unidades
soltas e uma dezena nova para levar. Agora as dezenas:
$2 + 1 + 1 = 4$ dezenas. Logo $27 + 15 = 42$.
\end{example}

\begin{omfigure}
\begin{tikzpicture}[scale=0.75]
  % 12 unidades: dez delas amarradas num feixe, duas soltas
  \foreach \i in {0,...,9} \draw[omThm, very thick]
    (0.24*\i,0) -- (0.24*\i,1.7);
  \draw[omExo, thick, rounded corners=3pt]
    (-0.24,0.62) rectangle (2.4,1.05);
  \node[below, font=\small] at (1.08,-0.25) {$10$ amarrados};
  \foreach \i in {0,1} \draw[omDef, very thick]
    (4.0+0.4*\i,0) -- (4.0+0.4*\i,1.7);
  \node[below, font=\small] at (4.2,-0.25) {sobram $2$};
  \node[right, font=\small, align=left] at (5.5,0.85)
    {$7 + 5 = 12$: uma dezena e $2$ unidades.\\
     O pacote é o vai-um.};
\end{tikzpicture}

{\small O vai-um é um pacote: dez unidades soltas viram uma dezena,
que passa para a coluna das dezenas.}
\end{omfigure}

\section{A conta armada}

\begin{method}[Adição em coluna]\label{met:g2:addition:column}
\begin{enumerate}
  \item Escreva os números um embaixo do outro, unidades sob
        unidades, dezenas sob dezenas;
  \item some as unidades; se o resultado passar de $9$, escreva o
        algarismo das unidades e leve um $1$ pequeno para cima das
        dezenas;
  \item some as dezenas (com o vai-um) e depois as centenas;
  \item confira com uma estimativa: arredonde cada número e some por
        alto.
\end{enumerate}
\end{method}

\begin{example}\label{ex:g2:addition:worked}
Calcule $268 + 156$:
\[
\begin{array}{r}
{\scriptstyle 1}\ {\scriptstyle 1}\ \phantom{0} \\[-3pt]
2\,6\,8 \\
+\ 1\,5\,6 \\
\hline
4\,2\,4
\end{array}
\]
Unidades: $8 + 6 = 14$: escreva $4$, vai $1$. Dezenas:
$6 + 5 + 1 = 12$: escreva $2$, vai $1$. Centenas:
$2 + 1 + 1 = 4$. Conferência por estimativa: mais ou menos
$270 + 160 = 430$, perto de $424$. \checkmark
\end{example}

\section{Somar de cabeça}

\begin{example}[Formar uma dezena]\label{ex:g2:addition:mental}
Em \omterm{def:g1:addition:def}{somas} pequenas, pule primeiro para um número redondo:
\[
48 + 6 = 48 + 2 + 4 = 50 + 4 = 54 .
\]
Some dezenas e unidades separadamente:
$35 + 23 = (30 + 20) + (5 + 3) = 50 + 8 = 58$. E um truque para
quem está quase na centena: $99 + 46 = 100 + 46 - 1 = 145$.
\end{example}

\begin{example}[Um probleminha]\label{ex:g2:addition:problem}
``A biblioteca tem $185$ livros; chegam $57$ novos. Com quantos
livros ela fica?'' Somando: $185 + 57 = 242$ (unidades
$5 + 7 = 12$, vai um\,\dots). A biblioteca fica com $242$ livros.
\end{example}

\section{Exercícios}

\begin{exercise}[$\star$]\label{exo:g2:addition:1}
Calcule (sem precisar do vai-um): $34 + 25$; \quad $126 + 43$; \quad
$507 + 82$.
\end{exercise}

\begin{exercise}[$\star$]\label{exo:g2:addition:2}
Calcule com vai-um: $47 + 8$; \quad $56 + 27$; \quad $65 + 19$.
\end{exercise}

\begin{exercise}[$\star$]\label{exo:g2:addition:3}
Calcule armando a conta: $268 + 145$; \quad $374 + 158$; \quad
$429 + 293$.
\end{exercise}

\begin{exercise}[$\star$]\label{exo:g2:addition:4}
Calcule de cabeça formando uma dezena: $38 + 5$; \quad $67 + 6$;
\quad $54 + 8$.
\end{exercise}

\begin{exercise}[$\star$]\label{exo:g2:addition:5}
Calcule somando dezenas e unidades separadamente: $42 + 36$; \quad
$53 + 24$; \quad $61 + 28$.
\end{exercise}

\begin{exercise}[$\star$]\label{exo:g2:addition:6}
Estime primeiro (arredonde para a dezena ou a centena mais próxima)
e depois calcule exatamente: $58 + 34$; \quad $296 + 187$.
\end{exercise}

\begin{exercise}[$\star$]\label{exo:g2:addition:7}
``Ana tem $128$ bolinhas e Tomás tem $95$. Quantas têm juntos?''
Escreva a adição e a frase de resposta.
\end{exercise}

\begin{exercise}[$\star$]\label{exo:g2:addition:8}
Complete o buraco nas unidades: $3\,?\, + 27 = 64$ (descubra o
algarismo das unidades que falta no primeiro número).
\end{exercise}

\begin{exercise}[$\star$]\label{exo:g2:addition:9}
Duas turmas vão ao espetáculo: $27$ e $28$ crianças, mais $6$
adultos. Quantas pessoas ao todo? (Some primeiro as duas turmas.)
\end{exercise}

\begin{exercise}[$\star\star$]\label{exo:g2:addition:10}
Encontre na lista $27$, $33$, $40$, $65$ três números cuja \omterm{def:g1:addition:def}{soma} seja
$100$. (Procure primeiro dois cujas unidades sejam amigas do dez.)
\end{exercise}
